
This paper reports an experiment (h-e1) designed to test whether contamination geometry strata can be formed from corpus-side signals (max 13-gram count and SBERT cosine similarity) and whether detector routing across these strata is feasible. The experiment does not produce routing results. Instead, it encounters a structural failure mode: random corpus streaming produces degenerate SBERT cosine distributions at 50,000-document scale, collapsing all 25,403 benchmark items into the lexical stratum. This failure mode, termed stratum collapse, is identified as a predictable consequence of using random corpus sampling for semantic similarity computation, and top-k nearest-neighbor retrieval is identified as a necessary prerequisite for semantic stratum formation.

The experiment additionally establishes two reproducible empirical findings: MMLU achieves n-gram recall = 1.0 against The Pile, C4, and FineWeb at 50K-document sampling; GSM8K achieves recall = 0.0 against the same corpora. These findings indicate that benchmark heterogeneity creates structurally different corpus overlap profiles that require benchmark-type-specific detection criteria.

Four implementation bugs are identified and documented, providing a concrete guide for future implementations of the detector families evaluated.

The routing hypothesis H-GeomRoute-v1 remains unresolved. The gate failure is attributed to experimental design and implementation failures, not to falsifying evidence. The next step is a Phase 0 redesign of h-e1 incorporating: (1) top-k FAISS nearest-neighbor retrieval per benchmark item for SBERT cosine computation, (2) corrected DCPDDDetector with two-model log-likelihood ratio, (3) benchmark-type-specific recall criteria, and (4) controlled synthetic contamination injection for stratum formation validation prior to testing on real corpus overlap.

